\section{Introduction}
A candidate-scoring interface takes a sequence, a task description, and a finite set of textual labels, then returns a distribution over those labels. Such an interface can expose several classification tasks through one model without generating free-form answers. Whether it remains accurate and reliable after adaptation to biological sequences is an empirical question: a model can exploit fixed templates, respond to candidate positions, or compress sequences without improving their representation.

We study this question using the open Laya typed-decision encoder \citep{laya2026}. Our scope is two supervised, closed-set tasks: DNA promoter detection and coarse protein structural classification. We compare raw sequence tokenization with a historical biological BPE vocabulary, and compare candidate scoring with an implemented fixed-class readout and a trained text-only control. The biological vocabulary is a functional input option, not an assumed source of improvement. No additional neural continual pretraining (CPT) is performed.

The test results distinguish interface utility from tokenization gains. With biological BPE, candidate scoring exceeds the implemented fixed head by 1.40 percentage points on DNA accuracy and 7.91 points on protein accuracy, averaged over three seeds. Yet the raw-input candidate model is more accurate than its biological-BPE counterpart on both tasks. Development-set diagnostics further show that candidate reorderings can change individual predictions while leaving aggregate accuracy nearly unchanged.

This study contributes (i) a reproducible adaptation and fixed-checkpoint evaluation protocol, (ii) matched-supervision comparisons of two input representations and two output interfaces, and (iii) diagnostics separating sequence dependence, candidate-order stability, and probability calibration. These are empirical contributions within the specified data and controls. The experiments do not establish unseen-task understanding, superiority to specialized biological encoders, or a new biological foundation model. We first place the study relative to existing sequence--text methods, then describe the protocol and its results.
